% Generated draft: requires mathematical and source review.
\documentclass[11pt]{article}
\usepackage[margin=25mm]{geometry}
\usepackage{fontspec}
\setmainfont{Latin Modern Roman}
\usepackage{amsmath,amssymb,mathtools,unicode-math}
\setmathfont{Latin Modern Math}
\usepackage{enumitem,needspace,xurl,hyperref}
\setlength{\parindent}{0pt}
\setlength{\parskip}{5pt}
\setlength{\emergencystretch}{3em}
\newcommand{\N}{\mathbb{N}}
\newcommand{\Z}{\mathbb{Z}}
\newcommand{\Q}{\mathbb{Q}}
\newcommand{\R}{\mathbb{R}}
\newcommand{\C}{\mathbb{C}}
\newcommand{\F}{\mathbb{F}}
\newcommand{\A}{\mathbb{A}}
\newcommand{\PP}{\mathbb{P}}
\newcommand{\E}{\mathbb{E}}
\newcommand{\eps}{\varepsilon}
\newcommand{\dd}{\,\mathrm{d}}
\begin{document}
\section*{UnsolvedMath ID: 4200005}
\textbf{Source problem code:} catalogue entry
\subsection{Definitions and mathematical statement}
For a subset $A$ of a metric space, let $\dim_H A$ denote its Hausdorff dimension, defined by the threshold at which its Hausdorff $q$-measure becomes zero. The source asks for fractal caustics in two settings, with the following existence targets:
\[
\forall s\in[1,2)\ \exists\text{ a convex billiard and a caustic }C,
\qquad \dim_H C=s,
\]
and
\[
\forall s\in[1,\infty)\ \exists\text{ a Riemannian manifold }M,
\ \exists p\in M,\qquad \dim_H C_p=s,
\]
where $C_p$ is the caustic of the point $p$. The manifold dimension is not fixed in the second assertion. Both assertions include $s=1$; the billiard dimensions strictly between $1$ and $2$ are the specifically fractal part highlighted by the research summary.

For the usual planar Birkhoff billiard, trajectories are straight segments in a convex domain and reflect with equal incidence and reflection angles at its boundary. A classical caustic is preserved by tangency of successive segments. The record does not define tangency or caustics for fractal sets, nor specify boundary regularity or whether the caustic itself must be convex. This matters: the boundary of a planar convex body is one-dimensional, so the classical convex-curve convention cannot simply be imposed on the requested dimensions above $1$.

For a geodesic flow, one common meaning of the caustic of $p$ is the conjugate locus, the critical values of the exponential map $\exp_p$ on its domain, away from the zero tangent vector. A cut locus or a fixed-time wavefront singular set is a different object. The source does not specify which meaning or regularity/completeness convention is intended. Thus the quantified existence targets are preserved, but a self-contained precise theorem requires the source's generalized caustic definitions in both settings.
\subsection{Short English statement}
Can convex billiards realize caustics of every Hausdorff dimension from $1$ to below $2$, and can point caustics on Riemannian manifolds realize every dimension at least $1$? The intended notion of a fractal caustic is not defined in the record.
\subsection{Sources}
\begin{itemize}
\item Local UnsolvedMath registry, permanent ID 4200005, exact source code catalogue entry, ``Fractal caustics''; proposed by Oliver Knill, 2000. Registry status open, classification OPEN-TRIAGE.
\item Oliver Knill, \emph{Open problems in Hamiltonian dynamics} (2000), item 5, ``FRACTAL CAUSTICS'', \url{https://people.math.harvard.edu/~knill/seminars/intr/}. Supplied HTML extraction accessed July 29, 2026. The page presents the item as open; no later authoritative resolution check is recorded.
\item The supplied research summary says existence of billiard caustics of Hausdorff dimension strictly between $1$ and $2$ appears open and distinguishes known nonsmooth non-fractal caustics. No precise reference or definition for those claims is supplied, and the summary does not resolve the separate manifold question.
\item Provenance flags status and rights NEEDS\_REVIEW: public access was verified, but redistribution permission was not established.
\end{itemize}
\end{document}
